\section{An exact unequal-rail hardtanh norm}
\label{sec:atomic}

In this section consider the full family of transversals, indexed by
$s\in\mathbb F_2^{\ell}$ independently, before retaining the column
subcode. Write $m=\ell$ temporarily. Let $\mathcal C$ be the span of the code characters $\chi_{A_s}$.
Let $\mathcal D_{\rm odd}$ be the span of all odd-degree characters
whose supports are contained in some $A_s$.
The odd-degree restriction is essential: a complete rail transversal
without the carrier has even degree and would otherwise defeat the
observability argument used later.

\begin{theorem}[Sharp atom bound and support stability]
\label{thm:atomic}
Suppose $m\ge2$, all rail sizes are even and lie in $[R,R+\Delta]$,
$\Delta\ge2$ is even, and $R\ge6\Delta$.
Every affine hardtanh atom $\phi_u(x)=\clip(u\cdot\bar x)$ satisfies
\[
 Q(\phi_u):=\sum_s w(n_s)|\widehat\phi_u(A_s)|\le\|u\|_2.
\]
Nonzero equality requires zero bias, zero irrelevant and unused-rail
weights, and magnitude one on the carrier and every coordinate of a
single selected transversal. Moreover
\begin{equation}\label{eq:atomic-stability}
 \|\phi_u-P_{\mathcal D_{\rm odd}}\phi_u\|_2^2
 \le2^8R^2\|u\|_2\bigl(\|u\|_2-Q(\phi_u)\bigr).
\end{equation}
Consequently $V(bg)=B_w(g)$ for every code score.
\end{theorem}

\subsection{Positive Fourier integrals and block balancing}

Allow an auxiliary threshold $v>0$ and write
$\clip_v(t)=\relu(t+v)-\relu(t-v)-v$.
A code coefficient of $\clip_v$ is twice the ReLU coefficient obtained
by adjoining a Boolean sign $\xi$ of weight $v$ and including it in
the support. Every augmented support has even degree. The even-degree
ReLU Fourier integral, up to its fixed degree sign, is
\begin{equation}\label{eq:relu-integral}
 \frac1\pi\int_0^\infty\cos(bt)
   \prod_{j\in A}\sin(u_jt)
   \prod_{j\notin A}\cos(u_jt)\,\frac{dt}{t^2}.
\end{equation}
Discard the absolute bias and irrelevant cosine factors, flip signs
of incoming weights, and take absolute values inside this integral.
Call the resulting weighted expression $H(u_0,\mathrm{rails};v)$,
where $u_0\ge0$ is the physical carrier magnitude. It bounds $Q$
and is homogeneous of degree one jointly in its arguments. The
auxiliary weight is excluded from the physical incoming norm.

H\"older on each even rail, applied to the disjoint union of
coefficient-index copies of the integration measure, bounds $H$ by
the geometric mean of its values on vectors obtained by choosing one
coordinate magnitude uniformly in each rail and repeating it throughout
that rail. The carrier and auxiliary weights stay fixed. This statement
concerns the positive absolute-integral expression, not an equality
between signed coefficients. Sampling preserves the expected squared
physical norm. Thus a norm bound for these balanced potentials
transfers to arbitrary coordinates by Jensen's inequality.

For balanced rails name their amplitudes $a_i\ge b_i\ge0$, with
respective sizes $p_i,q_i$, and put $D=\sum_i p_i$, $n=D+1$.
Let $J(c)$ be the weighted ReLU coefficient sum with two unpaired
shared coordinates, both of magnitude $c$, and the actual rail weights.
Every sine and cosine exponent in its integral is even, so the integral
has no sign cancellation. Cauchy--Schwarz on the two shared factors
in $H$ gives
\begin{equation}\label{eq:shared-cauchy}
 H(u_0,\mathrm{rails};v)\le2\sqrt{J(u_0)J(v)}.
\end{equation}

\subsection{A randomized minimum for unequal rails}

Write $\kappa_D^{\rm R}=2^{1-D}\binom{D-2}{(D-2)/2}$ for even $D$.
Consider a zero-bias ReLU ridge with selected rail coordinates $X_i$,
opposite coordinates $Y_i$, and any additional unpaired even block.
Let $c_s$ be its absolute coefficient on the transversal choosing $Y_i$
when $s_i=1$, together with the unpaired block.
Put $T_\varepsilon=\min_i(a_i+\varepsilon_i b_i)$.

\begin{lemma}[Unpaired shared block]\label{lem:random-min}
If the unpaired block consists of two coordinates of weight $c$, then
\[
 \sum_s c_s\le\kappa_{D+2}^{\rm R}\E\min(c,T_\varepsilon).
\]
If $q_i\ge p_i$, the coefficient mass using the opposite rail in
pair $i$ is at most
$\kappa_{D-p_i+q_i+2}^{\rm R}b_i\le\kappa_{D+2}^{\rm R}b_i$.
\end{lemma}
\begin{proof}
The fixed signs of the even-block integrals express the unweighted
sum as a top selected-parity coefficient after conditioning
\[
 \chi_{Y_i}=\eta_i\chi_{X_i},\qquad
 \eta_i=(-1)^{(q_i-p_i)/2},
\]
for every pair. To generate each conditional law, choose independent
uniform signs $\varepsilon_i$, set a distinguished opposite coordinate
$Y_{i,1}=\varepsilon_iX_{i,1}$, and condition the remaining opposite
coordinates on
$\chi_{Y_{i,\rm rest}}=\eta_i\varepsilon_i\chi_{X_{i,\rm rest}}$.
The distinguished selected coordinate now has nonnegative coefficient
$a_i+\varepsilon_i b_i$. The conditional laws do not depend on any
of these distinguished selected coordinates.

The normalized finite derivative of ReLU in a distinguished coordinate
is a nondecreasing function into $[0,1]$. It remains nondecreasing in
all other selected coordinates under these conditional laws: changing
an ordinary selected bit increases its direct contribution by $2a_j$
and reverses one opposite parity condition. Couple the two uniform
parity classes by flipping one fixed remaining opposite bit, changing
the total argument by $2a_j-2b_jY_{j,k}\ge0$.
An unpaired shared bit simply increases the argument by $2c$.

A nondecreasing function on $k$ Boolean coordinates with values in
$[0,1]$ has top Fourier coefficient at most
$2^{-k}\binom{k-1}{\lfloor(k-1)/2\rfloor}$ in absolute value.
Indeed average over permutations and express the resulting increasing
Hamming-level sequence as a convex combination of level thresholds;
their alternating binomial sums give this bound.
Apply it after taking a finite derivative in any distinguished selected
bit or in either shared bit. At fixed $\varepsilon$ the top coefficient
is at most $\kappa_{D+2}^{\rm R}\min(c,T_\varepsilon)$. Averaging
proves the first assertion.

For the marginal assertion, discard the cosine factor of the unused
selected rail $X_i$ in the positive integral. Treat $Y_i$ as an unpaired
selected block, and take its distinguished derivative of coefficient
$b_i$. The same argument gives
$\kappa_{D-p_i+q_i+2}^{\rm R}b_i$. Finally
$\kappa_{r+2}^{\rm R}/\kappa_r^{\rm R}=(r-1)/r<1$.
\end{proof}

\subsection{The exact norm bound}

The hardtanh cost $w(n)$ is increasing and its odd forward slopes are
at most three. Indeed
\[
 \frac{w(n+2)}{w(n)}=(1+1/n)\sqrt{1+2/n}\le(1+1/n)^2.
\]
The elementary central-binomial bound $w(n)\le2n$ therefore gives
$w(n+2)-w(n)\le4+2/n\le6$.
For completeness, the bound follows by induction from
$\kappa_{2r+1}\ge1/(2\sqrt r)$ for $r\ge1$, using the recurrence
$\kappa_{2r+3}/\kappa_{2r+1}=(2r+1)/(2r+2)$; $n=1$ is immediate.

Set $t=\E T_\varepsilon$ and $E_a=\sum_i(a_i-t)$. Monotonicity and
the slope bound imply, even when some opposite rails are shorter,
\[
 w(n_s)\le w(n)+3\sum_{i:q_i>p_i}(q_i-p_i)s_i.
\]
Lemma~\ref{lem:random-min} consequently yields
\[
 J(c)\le\frac{\sqrt n}{2}\E\min(c,T_\varepsilon)
       +3\kappa_{D+2}^{\rm R}\Delta\sum_i b_i.
\]
Cauchy--Schwarz, arithmetic--geometric mean, and concavity of the
minimum give
\[
 H\le J(u_0)+J(v)
 \le\frac{\sqrt n}{2}[\min(u_0,t)+t]
       +6\kappa_{D+2}^{\rm R}\Delta\sum_i b_i.
\]
For every $u_0,t\ge0$,
\[
 \frac n2[\min(u_0,t)+t]\le Dt+u_0.
\]
If $u_0\ge t$ this is immediate; otherwise the difference equals
$(D-1)(t-u_0)/2$. Also, for $i\ne j$,
\[
 t\le\frac{a_i+a_j-\max(|a_i-a_j|,b_i)}2,
\]
so summing over ordered pairs gives $\sum_i b_i\le2E_a$.
The recurrence shows $\kappa_n\sqrt n\le1$ for odd $n$; thus
$\kappa_{D+2}^{\rm R}\sqrt n\le1/2$. Combining these facts and
$R\ge6\Delta$ gives the stronger linear bound
\begin{equation}\label{eq:projection-bound}
 H\le\frac{Dt+u_0+6\Delta E_a}{\sqrt n}
 \le L:=\frac{\sum_i p_i a_i+u_0}{\sqrt n}
 \le\|u_{\rm physical}\|_2.
\end{equation}
The quantity $L$ is the projection onto the unit balanced ray on the
selected rails and physical carrier. If $u_*$ is that projection,
\begin{equation}\label{eq:balanced-distance}
 \|u-u_*\|_2^2=\|u\|_2^2-L^2
 \le2\|u\|_2(\|u\|_2-H).
\end{equation}
Block balancing and Jensen now prove the global atom inequality.
Equality forces zero unused weights and a common magnitude $a$ on
all selected physical coordinates and the carrier. Equality in the
shared-factor Cauchy inequality requires $\sin^2(at)$ proportional
to $\sin^2(vt)$ on an interval, hence $a=v$. At the standard threshold
$v=1$ the selected magnitudes are one. Bias and irrelevant weights
vanish because they only increase the full incoming norm.

\subsection{Polynomial stability into the odd downward space}

It suffices to control discarded support weights, rather than to force
all retained coordinates to be equal. By joint homogeneity, rescale the auxiliary threshold along with the
physical parameters to normalize the original relevant physical norm
to one. Let $V$ be its balanced rail sample. If
independent rail samples $X,Y$ have second moments $A\ge B$, then
\[
 \E\min(X^2,Y^2)\ge\frac{AB}{r_{\max}(A+B)}
                    \ge\frac B{2r_{\max}}.
\]
Choose deterministically the rail with larger second moment in each
pair. Its discarded physical squared norm is at most
$8R\E d_{\rm face}(V)^2$, since $R\le r_i^j\le2R$.
For $\epsilon=1-H(u)<1/2$, H\"older and $\E\|V\|^2=1$ imply
\[
 \E\left[1-\frac{H(V)}{\|V\|}\right]
 \le-\log H(u)\le2\epsilon,
 \qquad \|V\|^2\le r_{\max}\le2R.
\]
Using \eqref{eq:balanced-distance} bounds the discarded squared norm
by $64R^2\epsilon$. For $\epsilon\ge1/2$, use the trivial bound.
Bias and irrelevant squared norm are at most
$2\|u\|(\|u\|-Q(\phi_u))$, by the discarded-norm comparison.
Strip them and the chosen unused rails. The resulting zero-bias
hardtanh atom is odd and depends on one selected support, so lies in
$\mathcal D_{\rm odd}$. Hardtanh is one-Lipschitz and the augmented
Boolean coordinates have identity covariance. Parameter squared
distance therefore bounds the output $L^2$ residual, proving
\eqref{eq:atomic-stability} with the stated safe constant.
The exact variation norm follows from the atom inequality and the
finite parity-orbit representations.
